\subsection*{Chinese image histories improve spelling only under a supplied prefix}
To test whether familiar writing patterns could improve the next glyph, we indexed 8,424,544 consecutive glyph images from 192,230 other human answers. A control shuffled the same tiles within each answer. We kept the reader, action scores and set of available output images fixed. Given the correct preceding answer and the writing action, four-glyph histories recovered 431 of 1,525 exact missing cells. The neural writer alone recovered 162 and shuffled histories recovered 121. Ordered sequences gained 356 correct cells and lost 87 relative to the neural writer (Fig.~\ref{fig:bitmap-history}).

This advantage depended on the supplied history. Starting with an empty output, all three writers matched 25 complete references out of 571. None matched a reference in the 300 cases requiring input-absent ink. With the ordered writer, the first disagreement occurred before the first missing reference glyph in 228 of those cases and at that glyph in the remaining 72. Although exact references allow limited conclusions about alternative phrasings, the pixel writer had also failed the separate reading of supported responses. The local image sequences contain predictive evidence about Chinese spelling. They leave the initial choice of what to say unresolved.
